% Seminar 2 - Pagina 5. Fragment pentru % Seminar 2 - Pagina 5. Fragment pentru % Seminar 2 - Pagina 5. Fragment pentru % Seminar 2 - Pagina 5. Fragment pentru \input{seminars/pages/seminar2-pagina5-ro}.
% Fara preambul si fara intrerupere fortata de pagina.
% NOTE EDITORIALE:
% - Sub exercitiul 5 exista un calcul taiat (aparent g(x)-f(x)>0?);
%   nu a fost inclus in transcrierea curata.
% - In demonstratia exercitiului 5, pasul f(n_2) in A rezulta din
%   inegalitatea stricta f(n_2)<g(n_2); manuscrisul il noteaza prescurtat.
% - In exercitiul 7, conditia este citita f(n+1)>f(f(n)).
% - Ultimele randuri de pe pagina stabilesc numai f(n)>=n;
%   nu este transcrisa o continuare inexistenta a demonstratiei f=id.

\subsection*{Seminar 2 - Pagina 5}

\noindent\fbox{5}\quad Fie $f,g:\mathbb{N}\longrightarrow\mathbb{N}$,
$f$ injectivă, $g$ surjectivă. Dacă
\[
 f(n)\leq g(n)\qquad (\forall n\in\mathbb{N}),
\]
\noindent atunci $f=g$.

\medskip
\noindent\fbox{6}\quad Fie $f,g:\mathbb{N}\longrightarrow\mathbb{N}$. Dacă
\[
 \max(f,g)\ \text{este surjectivă},\qquad
 \min(f,g)\ \text{este injectivă},
\]
\noindent atunci $f=g$.

\medskip
\noindent\underline{Sol. 5.} Presupunem $f\neq g$. Atunci există
$n\in\mathbb{N}$ astfel încât
\[
 f(n)\neq g(n)\quad\Longrightarrow\quad f(n)<g(n).
\]
\noindent Fie
\[
 A=\{f(n)\mid f(n)\neq g(n)\}\neq\varnothing.
\]
\noindent $A$ este o mulţime de numere naturale; există
\[
 m=\min A=f(n_1)\in\mathbb{N}.
\]
\noindent Deoarece $g$ este surjectivă, există $n_2\in\mathbb{N}$ astfel încât
\[
 g(n_2)=m=f(n_1),\qquad n_1\neq n_2.
\]
\noindent Prin injectivitatea lui $f$,
\[
 f(n_1)\neq f(n_2).
\]
\noindent Dar
\[
 f(n_2)\leq g(n_2)=f(n_1)
 \quad\Longrightarrow\quad f(n_2)<f(n_1).
\]
\noindent Rezultă $f(n_2)\in A$ şi
\[
 f(n_2)<\min A,\qquad \text{fals}.
\]
\noindent Deci $A=\varnothing$, de unde $f=g$.

\medskip
\noindent\underline{Sol. 6.}
\[
 \min(f,g)\leq\max(f,g)
 \quad\Longrightarrow\quad
 \min(f,g)=\max(f,g)
 \quad\Longrightarrow\quad f=g.
\]

\bigskip
\noindent\fbox{7}\quad Fie
$f:\mathbb{N}^{*}\longrightarrow\mathbb{N}^{*}$ astfel încât
\[
 f(n+1)>f(f(n))\qquad (\forall n\in\mathbb{N}^{*}).
\]
\noindent Atunci $f=1_{\mathbb{N}^{*}}$.

\medskip
\noindent\underline{Sol.} $f(k)\geq k$, pentru orice
$k\in\mathbb{N}^{*}$.

\noindent Mai general:
\[
 n\geq k\quad\Longrightarrow\quad f(n)\geq k.
\]
\noindent Inducţie după $k$.

\noindent $k=1$: nu avem nimic de arătat.

\noindent $k\longmapsto k+1$: dacă $n\geq k+1$, atunci
\[
 n-1\geq k\quad\Longrightarrow\quad f(n-1)\geq k.
\]
\noindent Prin ipoteza de inducţie şi condiţia din enunţ,
\[
 f(n)>f\bigl(f(n-1)\bigr)\geq k
 \quad\Longrightarrow\quad f(n)>k.
\]
% Demonstratia continua, daca exista, pe urmatoarea pagina a manuscrisului.
.
% Fara preambul si fara intrerupere fortata de pagina.
% NOTE EDITORIALE:
% - Sub exercitiul 5 exista un calcul taiat (aparent g(x)-f(x)>0?);
%   nu a fost inclus in transcrierea curata.
% - In demonstratia exercitiului 5, pasul f(n_2) in A rezulta din
%   inegalitatea stricta f(n_2)<g(n_2); manuscrisul il noteaza prescurtat.
% - In exercitiul 7, conditia este citita f(n+1)>f(f(n)).
% - Ultimele randuri de pe pagina stabilesc numai f(n)>=n;
%   nu este transcrisa o continuare inexistenta a demonstratiei f=id.

\subsection*{Seminar 2 - Pagina 5}

\noindent\fbox{5}\quad Fie $f,g:\mathbb{N}\longrightarrow\mathbb{N}$,
$f$ injectivă, $g$ surjectivă. Dacă
\[
 f(n)\leq g(n)\qquad (\forall n\in\mathbb{N}),
\]
\noindent atunci $f=g$.

\medskip
\noindent\fbox{6}\quad Fie $f,g:\mathbb{N}\longrightarrow\mathbb{N}$. Dacă
\[
 \max(f,g)\ \text{este surjectivă},\qquad
 \min(f,g)\ \text{este injectivă},
\]
\noindent atunci $f=g$.

\medskip
\noindent\underline{Sol. 5.} Presupunem $f\neq g$. Atunci există
$n\in\mathbb{N}$ astfel încât
\[
 f(n)\neq g(n)\quad\Longrightarrow\quad f(n)<g(n).
\]
\noindent Fie
\[
 A=\{f(n)\mid f(n)\neq g(n)\}\neq\varnothing.
\]
\noindent $A$ este o mulţime de numere naturale; există
\[
 m=\min A=f(n_1)\in\mathbb{N}.
\]
\noindent Deoarece $g$ este surjectivă, există $n_2\in\mathbb{N}$ astfel încât
\[
 g(n_2)=m=f(n_1),\qquad n_1\neq n_2.
\]
\noindent Prin injectivitatea lui $f$,
\[
 f(n_1)\neq f(n_2).
\]
\noindent Dar
\[
 f(n_2)\leq g(n_2)=f(n_1)
 \quad\Longrightarrow\quad f(n_2)<f(n_1).
\]
\noindent Rezultă $f(n_2)\in A$ şi
\[
 f(n_2)<\min A,\qquad \text{fals}.
\]
\noindent Deci $A=\varnothing$, de unde $f=g$.

\medskip
\noindent\underline{Sol. 6.}
\[
 \min(f,g)\leq\max(f,g)
 \quad\Longrightarrow\quad
 \min(f,g)=\max(f,g)
 \quad\Longrightarrow\quad f=g.
\]

\bigskip
\noindent\fbox{7}\quad Fie
$f:\mathbb{N}^{*}\longrightarrow\mathbb{N}^{*}$ astfel încât
\[
 f(n+1)>f(f(n))\qquad (\forall n\in\mathbb{N}^{*}).
\]
\noindent Atunci $f=1_{\mathbb{N}^{*}}$.

\medskip
\noindent\underline{Sol.} $f(k)\geq k$, pentru orice
$k\in\mathbb{N}^{*}$.

\noindent Mai general:
\[
 n\geq k\quad\Longrightarrow\quad f(n)\geq k.
\]
\noindent Inducţie după $k$.

\noindent $k=1$: nu avem nimic de arătat.

\noindent $k\longmapsto k+1$: dacă $n\geq k+1$, atunci
\[
 n-1\geq k\quad\Longrightarrow\quad f(n-1)\geq k.
\]
\noindent Prin ipoteza de inducţie şi condiţia din enunţ,
\[
 f(n)>f\bigl(f(n-1)\bigr)\geq k
 \quad\Longrightarrow\quad f(n)>k.
\]
% Demonstratia continua, daca exista, pe urmatoarea pagina a manuscrisului.
.
% Fara preambul si fara intrerupere fortata de pagina.
% NOTE EDITORIALE:
% - Sub exercitiul 5 exista un calcul taiat (aparent g(x)-f(x)>0?);
%   nu a fost inclus in transcrierea curata.
% - In demonstratia exercitiului 5, pasul f(n_2) in A rezulta din
%   inegalitatea stricta f(n_2)<g(n_2); manuscrisul il noteaza prescurtat.
% - In exercitiul 7, conditia este citita f(n+1)>f(f(n)).
% - Ultimele randuri de pe pagina stabilesc numai f(n)>=n;
%   nu este transcrisa o continuare inexistenta a demonstratiei f=id.

\subsection*{Seminar 2 - Pagina 5}

\noindent\fbox{5}\quad Fie $f,g:\mathbb{N}\longrightarrow\mathbb{N}$,
$f$ injectivă, $g$ surjectivă. Dacă
\[
 f(n)\leq g(n)\qquad (\forall n\in\mathbb{N}),
\]
\noindent atunci $f=g$.

\medskip
\noindent\fbox{6}\quad Fie $f,g:\mathbb{N}\longrightarrow\mathbb{N}$. Dacă
\[
 \max(f,g)\ \text{este surjectivă},\qquad
 \min(f,g)\ \text{este injectivă},
\]
\noindent atunci $f=g$.

\medskip
\noindent\underline{Sol. 5.} Presupunem $f\neq g$. Atunci există
$n\in\mathbb{N}$ astfel încât
\[
 f(n)\neq g(n)\quad\Longrightarrow\quad f(n)<g(n).
\]
\noindent Fie
\[
 A=\{f(n)\mid f(n)\neq g(n)\}\neq\varnothing.
\]
\noindent $A$ este o mulţime de numere naturale; există
\[
 m=\min A=f(n_1)\in\mathbb{N}.
\]
\noindent Deoarece $g$ este surjectivă, există $n_2\in\mathbb{N}$ astfel încât
\[
 g(n_2)=m=f(n_1),\qquad n_1\neq n_2.
\]
\noindent Prin injectivitatea lui $f$,
\[
 f(n_1)\neq f(n_2).
\]
\noindent Dar
\[
 f(n_2)\leq g(n_2)=f(n_1)
 \quad\Longrightarrow\quad f(n_2)<f(n_1).
\]
\noindent Rezultă $f(n_2)\in A$ şi
\[
 f(n_2)<\min A,\qquad \text{fals}.
\]
\noindent Deci $A=\varnothing$, de unde $f=g$.

\medskip
\noindent\underline{Sol. 6.}
\[
 \min(f,g)\leq\max(f,g)
 \quad\Longrightarrow\quad
 \min(f,g)=\max(f,g)
 \quad\Longrightarrow\quad f=g.
\]

\bigskip
\noindent\fbox{7}\quad Fie
$f:\mathbb{N}^{*}\longrightarrow\mathbb{N}^{*}$ astfel încât
\[
 f(n+1)>f(f(n))\qquad (\forall n\in\mathbb{N}^{*}).
\]
\noindent Atunci $f=1_{\mathbb{N}^{*}}$.

\medskip
\noindent\underline{Sol.} $f(k)\geq k$, pentru orice
$k\in\mathbb{N}^{*}$.

\noindent Mai general:
\[
 n\geq k\quad\Longrightarrow\quad f(n)\geq k.
\]
\noindent Inducţie după $k$.

\noindent $k=1$: nu avem nimic de arătat.

\noindent $k\longmapsto k+1$: dacă $n\geq k+1$, atunci
\[
 n-1\geq k\quad\Longrightarrow\quad f(n-1)\geq k.
\]
\noindent Prin ipoteza de inducţie şi condiţia din enunţ,
\[
 f(n)>f\bigl(f(n-1)\bigr)\geq k
 \quad\Longrightarrow\quad f(n)>k.
\]
% Demonstratia continua, daca exista, pe urmatoarea pagina a manuscrisului.
.
% Fara preambul si fara intrerupere fortata de pagina.
% NOTE EDITORIALE:
% - Sub exercitiul 5 exista un calcul taiat (aparent g(x)-f(x)>0?);
%   nu a fost inclus in transcrierea curata.
% - In demonstratia exercitiului 5, pasul f(n_2) in A rezulta din
%   inegalitatea stricta f(n_2)<g(n_2); manuscrisul il noteaza prescurtat.
% - In exercitiul 7, conditia este citita f(n+1)>f(f(n)).
% - Ultimele randuri de pe pagina stabilesc numai f(n)>=n;
%   nu este transcrisa o continuare inexistenta a demonstratiei f=id.

\subsection*{Seminar 2 - Pagina 5}

\noindent\fbox{5}\quad Fie $f,g:\mathbb{N}\longrightarrow\mathbb{N}$,
$f$ injectivă, $g$ surjectivă. Dacă
\[
 f(n)\leq g(n)\qquad (\forall n\in\mathbb{N}),
\]
\noindent atunci $f=g$.

\medskip
\noindent\fbox{6}\quad Fie $f,g:\mathbb{N}\longrightarrow\mathbb{N}$. Dacă
\[
 \max(f,g)\ \text{este surjectivă},\qquad
 \min(f,g)\ \text{este injectivă},
\]
\noindent atunci $f=g$.

\medskip
\noindent\underline{Sol. 5.} Presupunem $f\neq g$. Atunci există
$n\in\mathbb{N}$ astfel încât
\[
 f(n)\neq g(n)\quad\Longrightarrow\quad f(n)<g(n).
\]
\noindent Fie
\[
 A=\{f(n)\mid f(n)\neq g(n)\}\neq\varnothing.
\]
\noindent $A$ este o mulţime de numere naturale; există
\[
 m=\min A=f(n_1)\in\mathbb{N}.
\]
\noindent Deoarece $g$ este surjectivă, există $n_2\in\mathbb{N}$ astfel încât
\[
 g(n_2)=m=f(n_1),\qquad n_1\neq n_2.
\]
\noindent Prin injectivitatea lui $f$,
\[
 f(n_1)\neq f(n_2).
\]
\noindent Dar
\[
 f(n_2)\leq g(n_2)=f(n_1)
 \quad\Longrightarrow\quad f(n_2)<f(n_1).
\]
\noindent Rezultă $f(n_2)\in A$ şi
\[
 f(n_2)<\min A,\qquad \text{fals}.
\]
\noindent Deci $A=\varnothing$, de unde $f=g$.

\medskip
\noindent\underline{Sol. 6.}
\[
 \min(f,g)\leq\max(f,g)
 \quad\Longrightarrow\quad
 \min(f,g)=\max(f,g)
 \quad\Longrightarrow\quad f=g.
\]

\bigskip
\noindent\fbox{7}\quad Fie
$f:\mathbb{N}^{*}\longrightarrow\mathbb{N}^{*}$ astfel încât
\[
 f(n+1)>f(f(n))\qquad (\forall n\in\mathbb{N}^{*}).
\]
\noindent Atunci $f=1_{\mathbb{N}^{*}}$.

\medskip
\noindent\underline{Sol.} $f(k)\geq k$, pentru orice
$k\in\mathbb{N}^{*}$.

\noindent Mai general:
\[
 n\geq k\quad\Longrightarrow\quad f(n)\geq k.
\]
\noindent Inducţie după $k$.

\noindent $k=1$: nu avem nimic de arătat.

\noindent $k\longmapsto k+1$: dacă $n\geq k+1$, atunci
\[
 n-1\geq k\quad\Longrightarrow\quad f(n-1)\geq k.
\]
\noindent Prin ipoteza de inducţie şi condiţia din enunţ,
\[
 f(n)>f\bigl(f(n-1)\bigr)\geq k
 \quad\Longrightarrow\quad f(n)>k.
\]
% Demonstratia continua, daca exista, pe urmatoarea pagina a manuscrisului.
